% chapters/10/sub-cpmk.tex
\begin{subcpmkbox}
Bab ini menuntaskan \SUBCPMK{2.3} secara utuh. Ketiga indikatornya diperiksa,
dan buktinya berupa program pencarian jalur beserta visualisasinya.

\vspace{0.5em}
\textbf{\SUBCPMK{2.3} Menerapkan \istilah{pathfinding} berbasis grid dengan
algoritma A* untuk navigasi NPC} --- ranah C4, \CPMK{2}, \CPL{03}.

\vspace{0.5em}
{\setlength{\tabcolsep}{4pt}\footnotesize%
\begin{tabularx}{\textwidth}{@{}C{1.9cm} L{4.0cm} L{2.5cm} Y@{}}
\toprule
\headrow \textbf{Indikator} & \textbf{Rumusan} & \textbf{Dibahas pada} & \textbf{Bukti} \\
\midrule
\ind{a} & Membangun representasi grid/graf berbiaya dari peta berhalangan &
Bagian~\ref{sec:peta-graf} &
Kode pembacaan peta dan penentuan tetangga, beserta pemeriksaan otomatis yang
menunjukkan jalur tidak pernah melewati dinding. \\
\addlinespace
\ind{b} & Menganalisis pengaruh heuristik terhadap optimalitas dan kinerja
pencarian &
Bagian~\ref{sec:heuristik} &
Tabel perbandingan tiga heuristik pada peta yang sama, beserta penjelasan
mengapa biaya jalurnya sama sedangkan banyaknya petak yang diperiksa
berbeda. \\
\addlinespace
\ind{c} & Merekonstruksi dan memvisualisasikan jalur hasil A* &
Bagian~\ref{sec:rekonstruksi} &
Kode rekonstruksi jalur, tangkapan layar pencarian, dan catatan pengamatan
pada lembar kerja praktikum. \\
\bottomrule
\end{tabularx}}
\end{subcpmkbox}
